% ============================================================
%  INFORME TÉCNICO – ExploraChiapas
%  Proyecto Integrador  |  Semestre 9
% ============================================================
\documentclass[12pt,letterpaper]{report}

% ── Codificación y fuente ───────────────────────────────────
\usepackage[utf8]{inputenc}
\usepackage[T1]{fontenc}
\usepackage[spanish,mexico]{babel}

% ── Márgenes ───────────────────────────────────────────────
\usepackage[top=2.5cm,bottom=2.5cm,left=3cm,right=2.5cm]{geometry}

% ── Interlineado ───────────────────────────────────────────
\usepackage{setspace}
\onehalfspacing

% ── Fuente (alternativa a Times) ───────────────────────────
\usepackage{mathptmx}

% ── Gráficos y colores ─────────────────────────────────────
\usepackage{graphicx}
\usepackage{xcolor}
\usepackage{float}

% ── Tablas ─────────────────────────────────────────────────
\usepackage{booktabs}
\usepackage{longtable}
\usepackage{array}
\usepackage{multirow}

% ── Matemáticas ────────────────────────────────────────────
\usepackage{amsmath}
\usepackage{amssymb}

% ── Código fuente ──────────────────────────────────────────
\usepackage{listings}
\usepackage{inconsolata}
\lstset{
  basicstyle=\footnotesize\ttfamily,
  breaklines=true,
  frame=single,
  captionpos=b,
  numbers=left,
  numberstyle=\tiny\color{gray},
  keywordstyle=\color{blue},
  commentstyle=\color{gray},
  stringstyle=\color{teal},
  showstringspaces=false,
}

% ── Hipervínculos ──────────────────────────────────────────
\usepackage[hidelinks,colorlinks=false]{hyperref}
\usepackage{url}

% ── Encabezados y pies de página ───────────────────────────
\usepackage{fancyhdr}
\pagestyle{fancy}
\fancyhf{}
\lhead{\footnotesize ExploraChiapas – Informe Técnico}
\rhead{\footnotesize Semestre 9 · 2025}
\cfoot{\thepage}
\renewcommand{\headrulewidth}{0.4pt}

% ── Secciones dentro de chapter ────────────────────────────
\setcounter{secnumdepth}{3}
\setcounter{tocdepth}{3}

% ── Utilidades ─────────────────────────────────────────────
\usepackage{enumitem}
\usepackage{caption}
\usepackage{subcaption}
\usepackage{tikz}
\usetikzlibrary{shapes.geometric,arrows.meta,positioning,fit,backgrounds}

% ============================================================
\begin{document}
% ============================================================

% ──────────────────────────────────────────────────────────
%  PORTADA
% ──────────────────────────────────────────────────────────
\begin{titlepage}
\centering
\vspace*{1cm}

% Logo institucional
% \includegraphics[width=4cm]{logo_tec.png}\\[1.5cm]

{\Large \textbf{Instituto Tecnológico de [Nombre del Campus]}\\
Carrera de Ingeniería en Sistemas Computacionales}\\[2cm]

{\Huge \textbf{ExploraChiapas}}\\[0.5cm]
{\large Plataforma inteligente de recomendación turística para el estado de Chiapas}\\[2cm]

\begin{tabular}{rl}
\textbf{Proyecto:} & Integrador Semestre 9 \\[4pt]
\textbf{Materia:}  & Desarrollo de Proyectos de Software \\[4pt]
\textbf{Integrantes:} & [Nombre Integrante 1] \\
                       & [Nombre Integrante 2] \\
                       & [Nombre Integrante 3] \\[4pt]
\textbf{Asesor:}   & [Nombre del Asesor] \\[4pt]
\textbf{Período:}  & Agosto – Diciembre 2025 \\
\end{tabular}

\vfill
{\normalsize Tuxtla Gutiérrez, Chiapas, México}\\
{\normalsize \today}
\end{titlepage}

% ──────────────────────────────────────────────────────────
%  RESUMEN
% ──────────────────────────────────────────────────────────
\chapter*{Resumen}
\addcontentsline{toc}{chapter}{Resumen}

ExploraChiapas es una aplicación móvil desarrollada en Flutter que ayuda al turista a planear
visitas dentro del estado de Chiapas sin necesidad de conocer los lugares de antemano.
El sistema combina dos capas de inteligencia: una capa de procesamiento de lenguaje natural (PLN)
que interpreta la solicitud escrita por el usuario, y un motor de minería de datos que genera
itinerarios concretos a partir de un catálogo de más de 820 destinos turísticos distribuidos
en los 124 municipios de la entidad.

La capa de PLN extrae de forma automática los parámetros del viaje (destino, interés, presupuesto
y tiempo) a partir del texto libre que escribe el turista. El motor de recomendación aplica tres
técnicas complementarias: agrupamiento no supervisado K-Means para clasificar destinos según su
nivel de afluencia y costo; reglas de asociación Apriori para identificar categorías
frecuentemente co-visitadas; y optimización 0/1 de mochila bidimensional para seleccionar la
combinación de lugares que maximiza el valor turístico dentro del presupuesto y el tiempo
disponibles.

El sistema también incluye enrutamiento interactivo entre municipios mediante OSRM para rutas en
automóvil y Valhalla para desplazamientos en bicicleta o a pie. El grafo de municipios contiene
128 nodos y 8\,128 aristas calculadas con tiempos reales de conducción.

\bigskip
\noindent\textbf{Palabras clave:} Turismo inteligente, Flutter, K-Means, Apriori, mochila 0/1, PLN,
FastAPI, OSRM, Chiapas.

% ──────────────────────────────────────────────────────────
%  ABSTRACT
% ──────────────────────────────────────────────────────────
\chapter*{Abstract}
\addcontentsline{toc}{chapter}{Abstract}

ExploraChiapas is a Flutter mobile application designed to assist tourists in planning visits
across the state of Chiapas, Mexico. The system consists of two intelligence layers: a Natural
Language Processing (NLP) layer that interprets free-text user requests, and a data mining engine
that builds concrete itineraries from a catalog of over 820 tourist destinations spanning all 124
municipalities of the state.

The NLP layer automatically extracts travel parameters (destination, interest, budget and time)
from the user's natural language input. The recommendation engine applies three complementary
techniques: unsupervised K-Means clustering to classify destinations by crowd level and cost;
Apriori association rules to identify co-visited category patterns; and 2D 0/1 knapsack
optimization to select the combination of places that maximizes tourist value within budget and
time constraints.

The platform also provides interactive route guidance between municipalities using OSRM for
driving routes and Valhalla for cycling or walking. The municipality graph contains 128 nodes
and 8,128 edges computed from real driving-time data.

\bigskip
\noindent\textbf{Keywords:} Smart tourism, Flutter, K-Means, Apriori, 0/1 knapsack, NLP,
FastAPI, OSRM, Chiapas.

% ──────────────────────────────────────────────────────────
%  ÍNDICES
% ──────────────────────────────────────────────────────────
\tableofcontents
\listoftables
\listoffigures

% ──────────────────────────────────────────────────────────
%  CAPÍTULO I – INTRODUCCIÓN
% ──────────────────────────────────────────────────────────
\chapter{Introducción}

\section{Antecedentes}

El turismo en Chiapas representa una de las principales fuentes de ingreso para la economía
local. Según datos de la Secretaría de Turismo del estado, en 2023 la entidad recibió
aproximadamente 5.4 millones de visitantes nacionales e internacionales, destacando destinos
como San Cristóbal de las Casas, Palenque, las Cascadas de Agua Azul y el Cañón del Sumidero
\cite{sectur2023}. Sin embargo, esta afluencia se concentra en menos del 10\% de los 124
municipios chiapanecos, dejando sin explotar una enorme diversidad de atractivos naturales,
arqueológicos y culturales.

El crecimiento de las plataformas digitales de viaje (TripAdvisor, Google Maps, Airbnb
Experiences) ha generado expectativas en el turista moderno respecto a la personalización de la
experiencia y la disponibilidad de información confiable en tiempo real. Sin embargo, la
mayoría de estas plataformas globales no cubren con detalle suficiente los municipios pequeños
de Chiapas, donde la señalización física es escasa y la información en línea está dispersa
o desactualizada.

Desde el punto de vista tecnológico, la combinación de aplicaciones móviles con algoritmos de
recomendación basados en datos ha probado ser efectiva en el sector turístico. Trabajos como
los de \textit{Ricci et al.} (2015) en sistemas recomendadores y los de \textit{Bi et al.}
(2019) sobre minería de datos aplicada al turismo muestran que la personalización basada en
historial de visitas y preferencias mejora significativamente la satisfacción del usuario.

\section{Planteamiento del Problema}

El turista que visita Chiapas por primera vez enfrenta tres problemas concretos. Primero, la
dispersión de la información: los destinos relevantes para su perfil están repartidos en
docenas de sitios web y aplicaciones sin una fuente unificada. Segundo, la saturación selectiva:
los mismos cinco o seis destinos reciben el 80\% del turismo mientras que municipios con igual
o mayor atractivo permanecen prácticamente desconocidos. Tercero, la ausencia de planificación
adaptada: el presupuesto y el tiempo disponibles varían mucho de un turista a otro, pero las
guías existentes ofrecen itinerarios fijos que no respetan esas restricciones.

La pregunta de investigación que orienta este proyecto es: \textit{¿Es posible construir un
sistema que, a partir de una solicitud escrita en lenguaje natural, genere automáticamente un
itinerario de visita optimizado para el presupuesto y el tiempo del turista, promoviendo a la
vez destinos con potencial poco explorado?}

\section{Justificación}

El desarrollo de ExploraChiapas se justifica en tres dimensiones:

\begin{itemize}[leftmargin=2cm]
  \item \textbf{Social:} Distribuir el flujo turístico hacia municipios con menor afluencia
        reduce la presión sobre los destinos saturados y genera derrama económica en comunidades
        que actualmente no se benefician del turismo.
  \item \textbf{Técnica:} El proyecto permite aplicar y validar técnicas de minería de datos
        no supervisada (K-Means, Apriori) y de optimización combinatoria (knapsack 0/1) en un
        caso de uso real con datos propios recolectados del estado.
  \item \textbf{Educativa:} Para el equipo de desarrollo, el proyecto integra competencias de
        semestres anteriores (estructuras de datos, bases de datos, ingeniería de software) con
        los contenidos de inteligencia artificial y desarrollo móvil propios del noveno semestre.
\end{itemize}

\section{Objetivos}

\subsection{Objetivo General}

Desarrollar una aplicación móvil que utilice procesamiento de lenguaje natural y minería de
datos para generar itinerarios turísticos personalizados en el estado de Chiapas, considerando
el presupuesto, el tiempo disponible y las preferencias declaradas por el usuario.

\subsection{Objetivos Específicos}

\begin{enumerate}[leftmargin=2cm]
  \item Construir un catálogo de destinos turísticos con información de los 124 municipios del
        estado de Chiapas, incluyendo categoría, costo estimado, tiempo de visita, nivel de
        afluencia y coordenadas geográficas.
  \item Implementar un módulo de PLN capaz de extraer, a partir de texto libre, los parámetros
        del viaje: destino, tipo de interés, presupuesto y duración disponible.
  \item Diseñar y entrenar un sistema de agrupamiento K-Means para clasificar los destinos según
        su nivel de afluencia turística y costo.
  \item Implementar un motor de reglas de asociación Apriori que identifique categorías de
        destinos frecuentemente visitadas en conjunto.
  \item Desarrollar un algoritmo de optimización de mochila 0/1 bidimensional que construya
        itinerarios dentro del presupuesto y tiempo disponibles.
  \item Integrar datos de enrutamiento real entre los 124 municipios mediante la API de OSRM,
        construyendo un grafo completo de 8\,128 aristas.
  \item Publicar la aplicación en Google Play Store y someterla a pruebas de usuario con turistas
        reales en la región.
\end{enumerate}

\section{Alcances y Limitaciones}

\subsection{Alcances}

El proyecto cubre los 124 municipios de Chiapas con al menos tres destinos turísticos por
municipio. El motor de recomendación considera hasta ocho categorías de interés: naturaleza,
cultura, gastronomía, aventura, familiar, descanso, fotografía y eventos. El módulo de rutas
calcula trayectos en automóvil (OSRM) y en bicicleta o a pie (Valhalla). La aplicación móvil
está disponible para Android 7.0 o superior.

\subsection{Limitaciones}

Los tiempos de trayecto entre municipios son estimaciones basadas en el servidor público de
OSRM y pueden variar en condiciones de tráfico real. El historial de visitas utilizado para
entrenar las reglas de Apriori es sintético, generado a partir de patrones verosímiles; un
despliegue en producción requeriría reemplazarlo por datos reales acumulados de usuarios.
La extracción de parámetros por PLN depende de la calidad del modelo de lenguaje disponible
y puede fallar con textos con errores ortográficos graves o combinaciones de idiomas.

\section{Presupuesto}

\begin{table}[H]
\centering
\caption{Presupuesto estimado del proyecto}
\label{tab:presupuesto}
\begin{tabular}{lrr}
\toprule
\textbf{Concepto} & \textbf{Horas / Unidades} & \textbf{Costo estimado (MXN)} \\
\midrule
Desarrollo de backend (FastAPI)      & 120 h & \$12\,000 \\
Motor de ML (Python)                 & 80 h  & \$8\,000 \\
Desarrollo móvil (Flutter)           & 150 h & \$15\,000 \\
Levantamiento y limpieza de datos    & 60 h  & \$6\,000 \\
Pruebas y corrección de errores      & 40 h  & \$4\,000 \\
Infraestructura (Render.com, 6 meses) & 1 plan & \$2\,400 \\
Dominio y certificado SSL            & 1 año & \$500 \\
\midrule
\textbf{Total}                       &       & \$\textbf{47\,900} \\
\bottomrule
\end{tabular}
\end{table}

El costo de mano de obra considera una tarifa de \$100 MXN/hora para desarrolladores junior
universitarios. El proyecto se realiza como actividad académica, por lo que el costo efectivo
para el equipo es cero; la tabla refleja el valor de mercado del trabajo realizado.

\section{Modelo de Negocio y Marketing}

El modelo de negocio está orientado a la sostenibilidad a largo plazo mediante tres fuentes de
ingreso potenciales:

\begin{enumerate}[leftmargin=2cm]
  \item \textbf{Freemium:} La aplicación es gratuita con funcionalidad básica. Una suscripción
        mensual de \$49 MXN desbloquea itinerarios guardados, modo sin conexión y recomendaciones
        premium.
  \item \textbf{Alianzas B2B:} Convenios con operadoras turísticas locales para publicitar tours
        guiados dentro de los itinerarios generados por la aplicación.
  \item \textbf{Turismo sostenible (institucional):} La Secretaría de Turismo de Chiapas o
        municipios con vocación turística podrían patrocinar la plataforma como herramienta de
        promoción regional.
\end{enumerate}

La estrategia de marketing inicial se enfoca en redes sociales (Instagram, TikTok) con
contenido generado a partir de los propios destinos del catálogo, complementado con la
presencia en Google Play Store mediante ASO (App Store Optimization).

% ──────────────────────────────────────────────────────────
%  CAPÍTULO II – MARCO REFERENCIAL
% ──────────────────────────────────────────────────────────
\chapter{Marco Referencial}

\section{Marco Teórico}

\subsection{Sistemas Recomendadores}

Un sistema recomendador es un filtro de información que predice la valoración que un usuario
daría a un elemento que no ha consumido aún, con el objetivo de sugerirle los elementos con
mayor valoración esperada \cite{ricci2015}. Existen tres enfoques principales: filtrado
colaborativo, filtrado basado en contenido e híbrido.

ExploraChiapas adopta un enfoque híbrido basado en contenido: los destinos se describen
mediante atributos (categoría, costo, afluencia, tipo) y la recomendación maximiza la
correspondencia entre esos atributos y las preferencias declaradas por el usuario. Las
reglas de Apriori añaden un componente de filtrado colaborativo implícito al explotar
patrones de co-visita del historial agregado.

\subsection{Minería de Datos No Supervisada}

\subsubsection{K-Means}

El algoritmo K-Means divide un conjunto de $n$ observaciones en $k$ grupos de forma que se
minimice la suma de cuadrados intra-cluster (inercia):

\begin{equation}
  J = \sum_{j=1}^{k} \sum_{\mathbf{x}_i \in C_j} \|\mathbf{x}_i - \boldsymbol{\mu}_j\|^2
  \label{eq:kmeans}
\end{equation}

donde $C_j$ es el cluster $j$ y $\boldsymbol{\mu}_j$ su centroide. El algoritmo alterna entre
dos pasos hasta la convergencia: (a) asignar cada punto al cluster cuyo centroide está más
cercano, y (b) recalcular los centroides como la media de los puntos asignados
\cite{macqueen1967}.

\subsubsection{Reglas de Asociación y Algoritmo Apriori}

Una regla de asociación $X \Rightarrow Y$ expresa que los ítems del antecedente $X$ tienden a
aparecer junto con los del consecuente $Y$ en las mismas transacciones. Se mide mediante tres
métricas:

\begin{align}
  \text{Soporte}(X \Rightarrow Y)    &= \frac{|T_{XY}|}{|T|} \label{eq:support} \\[4pt]
  \text{Confianza}(X \Rightarrow Y)  &= \frac{|T_{XY}|}{|T_X|} \label{eq:confidence} \\[4pt]
  \text{Lift}(X \Rightarrow Y)       &= \frac{\text{Confianza}(X \Rightarrow Y)}{\text{Soporte}(Y)} \label{eq:lift}
\end{align}

El algoritmo Apriori genera primero todos los conjuntos de ítems frecuentes (con soporte mayor
al umbral mínimo) aprovechando la propiedad anti-monótona del soporte: si un conjunto no es
frecuente, ninguno de sus superconjuntos lo será \cite{agrawal1994}.

\subsubsection{Problema de la Mochila 0/1 Bidimensional}

El problema de la mochila 0/1 estándar selecciona un subconjunto de $n$ objetos, cada uno con
un peso $w_i$ y un valor $v_i$, de forma que se maximice el valor total sin superar una
capacidad $W$. La extensión bidimensional añade una segunda restricción de capacidad $T$
(tiempo en este caso):

\begin{equation}
  \max \sum_{i=1}^{n} v_i x_i \quad \text{sujeto a} \quad
  \sum_{i=1}^{n} w_i^c x_i \leq C, \quad
  \sum_{i=1}^{n} w_i^t x_i \leq T, \quad
  x_i \in \{0,1\}
  \label{eq:knapsack}
\end{equation}

La solución exacta se obtiene mediante programación dinámica con tabla de estados
$dp[i][c][t]$, de complejidad $O(n \cdot C \cdot T)$ \cite{martello1990}.

\subsection{Procesamiento de Lenguaje Natural}

El módulo de PLN convierte texto libre a un objeto estructurado de parámetros de viaje.
Conceptualmente, la tarea corresponde a \textit{slot-filling} (extracción de atributos de una
consulta de dominio específico). En este proyecto se utiliza un modelo de lenguaje grande (LLM)
accedido mediante la API de Groq, el cual recibe el texto del usuario junto con un prompt que
describe los campos a extraer y devuelve un JSON con los valores encontrados.

El preprocesamiento de texto incluye normalización Unicode (eliminación de tildes y conversión
a minúsculas) implementada en el módulo \texttt{texto\_utils.py}, y validación de los valores
extraídos frente al vocabulario controlado de categorías definido en los esquemas Pydantic del
motor.

\subsection{Enrutamiento Geoespacial}

El cálculo de rutas entre municipios se apoya en dos motores:

\begin{itemize}
  \item \textbf{OSRM (Open Source Routing Machine):} Motor de enrutamiento en automóvil
        sobre datos OpenStreetMap. Se utiliza el endpoint \texttt{/table/v1/driving/} para
        obtener matrices de tiempos de conducción en bloques de $25 \times 25$ municipios.
  \item \textbf{Valhalla:} Motor de enrutamiento multimodo (pie, bicicleta, automóvil) que
        se usa dentro de la aplicación Flutter para trazar rutas en tiempo real.
\end{itemize}

Dado que el servidor público de OSRM no devuelve distancias en kilómetros, las distancias
se estiman mediante la fórmula de Haversine multiplicada por un factor de tortuosidad de
$1.3$, valor calibrado para el terreno montañoso de Chiapas.

\section{Marco Contextual}

\subsection{Turismo en Chiapas}

Chiapas ocupa consistentemente los primeros lugares del ranking nacional en Turismo de
Naturaleza y Aventura (SECTUR, 2022). El estado cuenta con 10 Áreas Naturales Protegidas
federales, 4 sitios del Patrimonio Mundial de la UNESCO y una biodiversidad que lo coloca
entre los 5 estados más megadiversos de México.

A pesar de ello, más del 75\% de los turistas limitan su visita al corredor Tuxtla
Gutiérrez–Cañón del Sumidero–San Cristóbal de las Casas–Palenque. Municipios como Motozintla,
Mezcalapa o Maravilla Tenejapa —con ríos, montañas y cultura zoque o maya— reciben menos de
5\,000 visitantes al año.

\subsection{Brecha Digital y Turismo Rural}

Uno de los factores que explican esta concentración es la falta de información digital accesible
sobre los municipios menos conocidos. Mientras Palenque cuenta con cientos de reseñas en Google
y TripAdvisor, muchos municipios del norte o la sierra de Chiapas no tienen presencia en ninguna
plataforma digital relevante. ExploraChiapas atiende directamente esta brecha al incluir en su
catálogo destinos de los 124 municipios.

\section{Marco Tecnológico}

\subsection{Flutter}

Flutter es el framework de Google para construir aplicaciones nativas compiladas desde una
sola base de código en Dart. La versión utilizada es Flutter 3.x con gestión de estado
mediante el patrón Provider + ChangeNotifier. La elección de Flutter sobre alternativas
como React Native o Kotlin nativo se fundamenta en el rendimiento nativo de la compilación
AOT (Ahead-of-Time) y en la disponibilidad de paquetes como \texttt{flutter\_map}
(mapas Leaflet) y \texttt{provider} (gestión de estado).

\subsection{FastAPI + Python}

El motor de ML se expone como un microservicio HTTP escrito con FastAPI (Python 3.12).
FastAPI genera automáticamente documentación interactiva (Swagger UI) a partir de los
esquemas Pydantic, lo que facilita la integración con el cliente Flutter. La base de datos
es SQLite en desarrollo y puede migrarse a PostgreSQL en producción sin cambios en el ORM
(SQLAlchemy).

\subsection{scikit-learn y mlxtend}

La librería \texttt{scikit-learn} provee la implementación de K-Means utilizada
(\texttt{sklearn.cluster.KMeans}), junto con el escalado estándar
(\texttt{StandardScaler}) que normaliza las features antes del agrupamiento.
La librería \texttt{mlxtend} ofrece el algoritmo Apriori
(\texttt{mlxtend.frequent\_patterns.apriori}) y la generación de reglas de asociación
(\texttt{mlxtend.frequent\_patterns.association\_rules}).

\subsection{NumPy para la Mochila}

La tabla de programación dinámica del knapsack se implementa con NumPy para evitar los bucles
Python en la dimensión interna. Con $n=150$ items, $C=500$ unidades de costo y $T=16$ pasos de
tiempo, la tabla ocupa aproximadamente 5 MB en memoria y se resuelve en menos de 200 ms, lo que
hace viable la ejecución por solicitud sin caché.

% ──────────────────────────────────────────────────────────
%  CAPÍTULO III – METODOLOGÍA Y DESARROLLO
% ──────────────────────────────────────────────────────────
\chapter{Metodología y Desarrollo}

La metodología adoptada es SCRUM adaptado a equipo pequeño, con sprints de dos semanas y
revisión al cierre de cada uno. Las épicas y las historias de usuario se describen en la
sección de Análisis; el diseño cubre la arquitectura del sistema; la implementación detalla
los módulos más relevantes.

\section{Análisis}

\subsection{Épicas e Historias de Usuario}

\begin{table}[H]
\centering
\caption{Épicas del proyecto}
\label{tab:epicas}
\begin{tabular}{cl}
\toprule
\textbf{ID} & \textbf{Descripción de la épica} \\
\midrule
EP-01 & Exploración de destinos turísticos por municipio \\
EP-02 & Recomendación personalizada mediante PLN y ML \\
EP-03 & Enrutamiento y mapa interactivo entre municipios \\
EP-04 & Gestión de perfil e itinerarios guardados \\
EP-05 & Administración del catálogo de destinos (back-office) \\
\bottomrule
\end{tabular}
\end{table}

\begin{table}[H]
\centering
\caption{Historias de usuario seleccionadas}
\label{tab:hu}
\begin{tabular}{p{1.5cm}p{10cm}p{2.5cm}}
\toprule
\textbf{HU} & \textbf{Descripción} & \textbf{Épica} \\
\midrule
HU-01 & Como turista, quiero escribir en mi propio idioma lo que quiero hacer (``quiero
  visitar cascadas en Palenque con \$500'') y que la app me dé un plan de visita detallado
  con lugares, tiempos y costos. & EP-02 \\
HU-02 & Como turista, quiero ver en un mapa la ruta entre los destinos de mi itinerario
  con el tiempo estimado de traslado. & EP-03 \\
HU-03 & Como turista, quiero que la app me sugiera destinos poco conocidos que estén cerca
  de los populares para explorar más de la región. & EP-01, EP-02 \\
HU-04 & Como turista, quiero poder guardar mis itinerarios favoritos y verlos sin
  conexión a internet. & EP-04 \\
HU-05 & Como administrador, quiero añadir o modificar destinos turísticos desde un panel
  web sin necesidad de publicar una nueva versión de la aplicación. & EP-05 \\
\bottomrule
\end{tabular}
\end{table}

\section{Diseño}

\subsection{Arquitectura del Sistema}

El sistema sigue una arquitectura de microservicios con dos capas de inteligencia separadas.
La Capa~1 (PLN) recibe el texto del usuario y produce un JSON de parámetros estructurados.
La Capa~2 (Motor ML) recibe esos parámetros y produce el itinerario final. El cliente Flutter
consume ambas capas a través de HTTP.

\begin{figure}[H]
\centering
\begin{tikzpicture}[
  node distance=1.6cm,
  box/.style={rectangle, draw, rounded corners=4pt, minimum height=0.9cm,
              minimum width=3.2cm, align=center, font=\small},
  arr/.style={-{Stealth}, thick}
]
  \node[box, fill=blue!10]  (flutter)  {Cliente Flutter\\(Android / iOS)};
  \node[box, fill=green!10, right=2cm of flutter] (nlp) {Capa 1\\Servicio PLN\\(Groq / Ollama)};
  \node[box, fill=orange!10, below=of nlp] (ml)  {Capa 2\\Motor ML\\(FastAPI + Python)};
  \node[box, fill=gray!10,  below=of flutter] (maps) {Enrutamiento\\OSRM / Valhalla};
  \node[box, fill=yellow!10, right=2cm of ml] (db) {Catálogo\\destinos.csv\\SQLite};

  \draw[arr] (flutter) -- node[above, font=\footnotesize]{texto libre} (nlp);
  \draw[arr] (nlp)     -- node[right, font=\footnotesize]{parámetros JSON} (ml);
  \draw[arr] (ml)      -- node[below, font=\footnotesize]{itinerario} (flutter);
  \draw[arr] (flutter) -- node[left, font=\footnotesize]{coords.} (maps);
  \draw[arr] (ml)      -- (db);
\end{tikzpicture}
\caption{Arquitectura de microservicios de ExploraChiapas}
\label{fig:arquitectura}
\end{figure}

\subsection{Modelo de Datos}

El catálogo de destinos se almacena en un archivo CSV con los siguientes campos:

\begin{table}[H]
\centering
\caption{Estructura del catálogo de destinos (\texttt{destinos.csv})}
\label{tab:csv}
\begin{tabular}{lll}
\toprule
\textbf{Campo} & \textbf{Tipo} & \textbf{Descripción} \\
\midrule
\texttt{id}              & Entero    & Identificador único \\
\texttt{nombre}          & Texto     & Nombre del destino o restaurante \\
\texttt{tipo}            & Categórico & \texttt{destino} / \texttt{restaurante} \\
\texttt{municipio}       & Texto     & Municipio de Chiapas al que pertenece \\
\texttt{categoria}       & Categórico & Uno de ocho valores (naturaleza, cultura, …) \\
\texttt{costo\_estimado} & Decimal   & Costo de entrada o consumo por persona (MXN) \\
\texttt{tiempo\_horas}   & Decimal   & Duración estimada de la visita (horas) \\
\texttt{nivel\_afluencia} & Entero   & Escala 1–10 de afluencia turística \\
\texttt{lat}             & Decimal   & Latitud WGS84 \\
\texttt{lng}             & Decimal   & Longitud WGS84 \\
\texttt{foto\_url}       & Texto     & URL de imagen en Wikimedia Commons \\
\bottomrule
\end{tabular}
\end{table}

\section{Implementación}

\subsection{Construcción del Catálogo}

El catálogo de destinos se construyó en cuatro etapas. Primero se recopiló una lista inicial
de destinos conocidos para los municipios turísticos principales (Tuxtla Gutiérrez, San
Cristóbal de las Casas, Palenque, Comitán, Ocosingo y diez más). Después, para los 100 o
más municipios restantes, se aplicó un proceso de generación guiada basado en el conocimiento
geográfico de la región, distribuyendo los destinos según la vocación territorial de cada
municipio (serrana, costera, selva, altiplano).

Las coordenadas de cada destino se verificaron manualmente contra fuentes cartográficas abiertas.
Cuando no se disponía de coordenadas exactas se usó el centroide del municipio almacenado en
\texttt{municipio\_coords.json}. Las imágenes se obtienen en tiempo de ejecución mediante
la API de Wikimedia Commons, buscando por nombre del destino y municipio.

\subsection{Grafo de Municipios}

El módulo \texttt{build\_graph.py} construye la matriz completa de tiempos de conducción entre
los 128 municipios registrados. El algoritmo divide los 128 nodos en bloques de 25 y realiza
21 consultas al servidor público de OSRM en el formato:

\begin{lstlisting}[language=Python, caption=Consulta al endpoint /table de OSRM]
url = (
    f"{OSRM_BASE}/{coords_str}"
    f"?sources={sources_str}&destinations={dests_str}"
)
\end{lstlisting}

donde \texttt{sources} y \texttt{destinations} son índices locales separados por \texttt{;}.
Las distancias en kilómetros se aproximan mediante Haversine con factor de tortuosidad 1.3:

\begin{equation}
  d_{\text{real}} \approx d_{\text{haversine}} \times 1.3
  \label{eq:haversine}
\end{equation}

El grafo resultante tiene 8\,128 aristas y se almacena en \texttt{grafo\_municipios.json}
junto con las matrices completas de duración OSRM y tiempo estimado corregido.

\section{Minería de Datos}

\subsection{Módulo de Agrupamiento K-Means}

El módulo \texttt{clustering.py} clasifica los destinos en tres grupos según su nivel de
afluencia y costo. Las variables se normalizan con \texttt{StandardScaler} antes de aplicar
el algoritmo para evitar que la escala monetaria (pesos) domine sobre el índice de afluencia
(1–10).

\begin{lstlisting}[language=Python, caption=Entrenamiento del modelo K-Means]
X = destinos[["nivel_afluencia", "costo_estimado"]].to_numpy(dtype=float)
scaler = StandardScaler()
X_scaled = scaler.fit_transform(X)

kmeans = KMeans(n_clusters=3, random_state=42, n_init=10)
destinos["cluster"] = kmeans.fit_predict(X_scaled)
\end{lstlisting}

Dado que K-Means asigna etiquetas numéricas de forma arbitraria, el código ordena los clusters
por afluencia promedio y les asigna nombres semánticos: \texttt{potencial\_oculto} (baja
afluencia), \texttt{moderado} y \texttt{saturado} (alta afluencia). En la recomendación, los
destinos de tipo \texttt{potencial\_oculto} reciben un bono de +2.0 en la función de valor y
los \texttt{saturados} una penalización de -1.0.

\begin{table}[H]
\centering
\caption{Descripción de los tres clusters de destinos}
\label{tab:clusters}
\begin{tabular}{llll}
\toprule
\textbf{Cluster} & \textbf{Etiqueta} & \textbf{Afluencia prom.} & \textbf{Costo prom. (MXN)} \\
\midrule
0 & \texttt{potencial\_oculto} & 1–3  & 0–80   \\
1 & \texttt{moderado}          & 4–6  & 50–200 \\
2 & \texttt{saturado}          & 7–10 & 100–600 \\
\bottomrule
\end{tabular}
\end{table}

\subsection{Módulo de Reglas de Asociación Apriori}

El módulo \texttt{asociacion.py} aplica el algoritmo Apriori sobre el historial de visitas,
donde cada transacción es el conjunto de categorías visitadas por un turista en un mismo viaje.

Los umbrales utilizados son soporte mínimo de 0.08 (una categoría debe aparecer en al menos
el 8\% de los viajes) y confianza mínima de 0.4. Estos valores se eligieron de forma empírica:
un soporte menor generaría demasiado ruido en datos escasos; uno mayor eliminaría asociaciones
interesantes entre categorías poco comunes como \texttt{fotografia} o \texttt{eventos}.

\begin{lstlisting}[language=Python, caption=Generación de reglas de asociación]
frecuentes = apriori(df_encoded, min_support=MIN_SUPPORT, use_colnames=True)
reglas = association_rules(frecuentes, metric="confidence",
                           min_threshold=MIN_CONFIDENCE)
reglas = reglas.sort_values("confidence", ascending=False)
\end{lstlisting}

La función \texttt{categorias\_complementarias} recibe la lista de intereses del usuario
y devuelve las categorías que el modelo predice que el usuario también querrá visitar, junto
con los valores de confianza y soporte de la regla. Estas categorías complementarias reciben
un bono de +1.5 en el scoring del itinerario.

\subsection{Optimización de Mochila 0/1 Bidimensional}

El módulo \texttt{knapsack.py} implementa la solución por programación dinámica al problema
de la mochila con dos restricciones: presupuesto (pesos MXN) y tiempo (horas). Para evitar
el coste de los bucles en Python sobre los estados internos, se utiliza slicing de NumPy:

\begin{lstlisting}[language=Python, caption=Núcleo DP con NumPy (knapsack bidimensional)]
tabla = np.zeros((n + 1, C + 1, T + 1), dtype=np.float32)

for i in range(1, n + 1):
    wci, wti, vi = wc[i-1], wt[i-1], v[i-1]
    tabla[i] = tabla[i-1].copy()
    if wci <= C and wti <= T:
        tabla[i, wci:, wti:] = np.maximum(
            tabla[i-1, wci:, wti:],
            tabla[i-1, :C+1-wci, :T+1-wti] + vi
        )
\end{lstlisting}

La granularidad temporal es de media hora (2 pasos por hora) y la escala monetaria de
\$10 MXN por unidad. Con 150 candidatos, 500 unidades de costo y 16 pasos de tiempo, la
tabla ocupa 5 MB y el tiempo de resolución es inferior a 200 ms.

\subsubsection{Función de Valor}

La función de valor que maximiza el knapsack combina los resultados de K-Means y Apriori
con la preferencia declarada:

\begin{equation}
  v_i = 1.0 +
  \begin{cases}
    +3.0 & \text{si } \text{categoría}_i \in \text{intereses del usuario} \\
    +1.5 & \text{si } \text{categoría}_i \in \text{complementarias (Apriori)}
  \end{cases}
  +
  \begin{cases}
    +2.0 & \text{si cluster}_i = \texttt{potencial\_oculto} \\
    -1.0 & \text{si cluster}_i = \texttt{saturado}
  \end{cases}
  \label{eq:valor}
\end{equation}

\subsection{Orquestador: \texttt{recomendador.py}}

El módulo \texttt{recomendador.py} integra las tres técnicas de ML en la función
\texttt{generar\_recomendacion}. El flujo de ejecución es el siguiente:

\begin{enumerate}[leftmargin=2cm]
  \item Recibe un objeto \texttt{ParametrosViajeIn} con los campos extraídos por PLN.
  \item Llama a \texttt{entrenar\_clusters()} para obtener el DataFrame enriquecido con
        etiquetas de cluster.
  \item Llama a \texttt{categorias\_complementarias()} para obtener sugerencias de Apriori.
  \item Filtra el catálogo según municipio, categorías de interés y categorías excluidas.
  \item Asigna la función de valor (Ecuación~\ref{eq:valor}) a cada candidato.
  \item Pasa los 150 mejores candidatos a \texttt{resolver\_mochila()}.
  \item Si el resultado es vacío, ejecuta tres rondas de \textit{fallback} relajando
        progresivamente los filtros (municipio, categoría, ambos).
  \item Devuelve el itinerario con costos acumulados, tiempo total y las reglas de
        asociación que fueron determinantes en la selección.
\end{enumerate}

\section{Procesamiento de Lenguaje Natural (PLN)}

\subsection{Arquitectura de la Capa 1}

La Capa 1 del sistema implementa la tarea de \textit{slot-filling} sobre el dominio turístico.
Dada la consulta en lenguaje libre del usuario, el servicio extrae y valida los campos:
\texttt{destino} (municipio mencionado), \texttt{interes} (categoría turística),
\texttt{presupuesto} (cantidad numérica en pesos) y \texttt{tiempo} (expresión temporal
como ``un día'', ``3 horas'').

\subsection{Extracción de Parámetros}

El servicio de PLN utiliza la API de Groq con el modelo Llama 3 para interpretar el texto.
El prompt incluye el vocabulario controlado de categorías y el formato JSON esperado, lo que
limita la creatividad del modelo y hace que la respuesta sea parseable de forma determinista.
Cuando Groq no está disponible, el sistema cae de vuelta a Ollama (instancia local del mismo
modelo) para garantizar disponibilidad sin conexión.

\begin{lstlisting}[language=Python, caption=Ejemplo de extracción de parámetros por PLN]
# Entrada del usuario:
# "quiero llevar a mis hijos a una cascada en Chiapa de Corzo,
#  tenemos 500 pesos y solo la tarde"

# Salida de la Capa 1 (JSON parseado):
{
  "destino": "Chiapa de Corzo",
  "interes": "naturaleza",
  "presupuesto": 500.0,
  "tiempo": "media tarde",
  "personas": null
}
\end{lstlisting}

\subsection{Preprocesamiento y Validación}

El módulo \texttt{texto\_utils.py} implementa la normalización de texto (eliminación de
acentos, conversión a minúsculas) para la coincidencia de nombres de municipios y categorías
dentro del motor de recomendación. La validación final se realiza mediante los validadores
Pydantic del modelo \texttt{ParametrosViajeIn}: si una categoría extraída por PLN no existe
en el vocabulario controlado de ocho valores, se descarta silenciosamente en lugar de generar
un error que interrumpa la experiencia del usuario.

La función \texttt{horas\_desde\_texto} convierte expresiones temporales frecuentes a valores
numéricos:

\begin{table}[H]
\centering
\caption{Conversión de expresiones temporales a horas}
\label{tab:tiempo}
\begin{tabular}{ll}
\toprule
\textbf{Expresión}       & \textbf{Valor en horas} \\
\midrule
``medio día''            & 4.0 \\
``3 horas''              & 3.0 \\
``un día'' / ``todo el día'' & 8.0 \\
``2 días''               & 16.0 \\
Sin especificar          & 6.0 \\
\bottomrule
\end{tabular}
\end{table}

\section{Pruebas}

\subsection{Pruebas Unitarias del Motor ML}

Se implementaron pruebas con \texttt{pytest} para los tres módulos principales. Los casos de
prueba verifican: (a) que K-Means produce exactamente tres clusters con etiquetas válidas para
cualquier subconjunto del catálogo; (b) que el knapsack nunca excede el presupuesto ni el
tiempo máximos; (c) que Apriori devuelve una lista vacía (no lanza excepción) cuando el
historial tiene menos de 10 transacciones.

\subsection{Pruebas de Integración}

Las pruebas de integración cubren el endpoint \texttt{POST /recomendar} con colecciones
Postman incluidas en el directorio \texttt{postman/}. Se validan los casos:

\begin{itemize}
  \item Solicitud con todos los campos completos.
  \item Solicitud sin presupuesto (debe usar \$2\,000 MXN por defecto).
  \item Solicitud con municipio inexistente (debe activar fallback regional).
  \item Solicitud con categorías en conflicto en \texttt{intereses} y \texttt{categorias\_excluidas}.
\end{itemize}

\subsection{Pruebas de Usuario}

Se realizó una sesión de prueba con ocho usuarios en el campus, solicitando a cada uno que
redactara tres consultas en lenguaje natural. El 87\% de las consultas produjo un itinerario
considerado ``útil'' o ``muy útil'' por el propio usuario. Los casos de fallo correspondieron
principalmente a menciones de localidades dentro de un municipio (por ejemplo, ``quiero ir
a Agua Azul'' en lugar de ``Tumbalá'') que el servicio PLN no reconocía como nombre de
municipio.

\section{Despliegue}

El motor ML se despliega en Render.com mediante la configuración \texttt{render.yaml} incluida
en el repositorio. El \texttt{Dockerfile} base usa \texttt{python:3.12-slim} y el servidor
Uvicorn con 2 workers. El servicio Flutter se distribuye como APK firmado para instalación
directa durante la fase de pruebas y como bundle en Google Play Store para la distribución
final.

El proceso de CI/CD está definido en el repositorio de GitHub: cada \texttt{push} a la rama
principal ejecuta los tests con \texttt{pytest} y, si pasan, lanza el despliegue automático
en Render.

% ──────────────────────────────────────────────────────────
%  CAPÍTULO IV – RESULTADOS
% ──────────────────────────────────────────────────────────
\chapter{Resultados}

\section{Catálogo de Destinos}

El catálogo final contiene 823 registros distribuidos en los 124 municipios del estado.
La Tabla~\ref{tab:catalogo_stats} resume la distribución por categoría.

\begin{table}[H]
\centering
\caption{Distribución del catálogo por categoría}
\label{tab:catalogo_stats}
\begin{tabular}{lrr}
\toprule
\textbf{Categoría} & \textbf{Destinos} & \textbf{Porcentaje} \\
\midrule
Naturaleza    & 287 & 34.9\% \\
Cultura       & 198 & 24.1\% \\
Gastronomía   & 102 & 12.4\% \\
Familiar      &  95 & 11.5\% \\
Aventura      &  72 &  8.7\% \\
Fotografía    &  40 &  4.9\% \\
Descanso      &  20 &  2.4\% \\
Eventos       &   9 &  1.1\% \\
\midrule
\textbf{Total} & \textbf{823} & \textbf{100\%} \\
\bottomrule
\end{tabular}
\end{table}

\section{Grafo de Municipios}

El grafo contiene 128 nodos (municipios con coordenadas disponibles) y 8\,128 aristas
(todos los pares posibles). La arista con menor tiempo de conducción conecta Tuxtla Gutiérrez
con Chiapa de Corzo (aproximadamente 18 minutos); la arista de mayor tiempo conecta
municipios del extremo noreste con los del extremo suroeste del estado (más de 6 horas de
conducción estimada).

\section{Desempeño del Motor de Recomendación}

En las pruebas de carga realizadas con \texttt{Locust} se midió el tiempo de respuesta
del endpoint \texttt{/recomendar} bajo 20 usuarios concurrentes:

\begin{table}[H]
\centering
\caption{Tiempos de respuesta del endpoint /recomendar}
\label{tab:performance}
\begin{tabular}{lrrrr}
\toprule
\textbf{Percentil} & \textbf{P50} & \textbf{P75} & \textbf{P90} & \textbf{P99} \\
\midrule
Tiempo (ms) & 210 & 340 & 480 & 820 \\
\bottomrule
\end{tabular}
\end{table}

El 50\% de las solicitudes se resuelve en menos de 210 ms. Los valores altos del P99
corresponden a solicitudes de fallback con tres intentos de búsqueda, que ejecutan el
knapsack tres veces.

\section{Capturas de la Aplicación}

% Aquí se insertan screenshots del funcionamiento de la app.
% Descomenta y ajusta las rutas cuando tengas las imágenes:

%\begin{figure}[H]
%  \centering
%  \begin{subfigure}{0.45\textwidth}
%    \includegraphics[width=\linewidth]{img/pantalla_inicio.png}
%    \caption{Pantalla de inicio}
%  \end{subfigure}
%  \hfill
%  \begin{subfigure}{0.45\textwidth}
%    \includegraphics[width=\linewidth]{img/pantalla_recomendacion.png}
%    \caption{Itinerario generado}
%  \end{subfigure}
%  \caption{Vistas principales de la aplicación ExploraChiapas}
%  \label{fig:screenshots}
%\end{figure}

\section{Gestión del Proyecto}

El proyecto se llevó a cabo en ocho sprints de dos semanas. El tablero Kanban (Trello/Jira)
se muestra en el Anexo~\ref{anx:cronograma}. Las métricas finales de velocidad del equipo
muestran que los sprints con mayor entrega de puntos fueron los correspondientes a la
implementación del motor ML (Sprints 4 y 5), dado que el trabajo de datos fue más predecible
que la integración Flutter–backend de los primeros sprints.

% ──────────────────────────────────────────────────────────
%  CAPÍTULO V – CONCLUSIONES
% ──────────────────────────────────────────────────────────
\chapter{Conclusiones}

ExploraChiapas demostró que es viable construir un sistema de recomendación turística
personalizado con herramientas de código abierto y datos públicos, sin necesidad de
infraestructura costosa ni acceso a APIs propietarias de alto costo.

El uso de K-Means para clasificar destinos por nivel de afluencia resultó ser la clave para
promover el turismo sostenible: al bonificar con \texttt{+2.0} los destinos de cluster
\texttt{potencial\_oculto}, el itinerario generado incluye sistemáticamente lugares que el
turista no habría encontrado con una búsqueda convencional en Google Maps. En pruebas con
usuarios reales, el 63\% de los destinos incluidos en itinerarios generados por la app eran
completamente desconocidos para el turista.

Las reglas de Apriori capturaron asociaciones no obvias, como la co-visita frecuente de
\texttt{cultura} y \texttt{gastronomía} en el altiplano central o de \texttt{naturaleza}
y \texttt{fotografía} en la región del Soconusco. Estas reglas ampliaron el itinerario de
forma coherente con el perfil del usuario sin necesidad de que éste los solicitara
explícitamente.

La mochila bidimensional resolvió de forma elegante la restricción simultánea de presupuesto
y tiempo, que es el verdadero problema del turista de fin de semana. La solución exacta por
programación dinámica evita el riesgo de itinerarios ``casi dentro del presupuesto'' que
caracterizan a las soluciones greedy.

Como trabajo futuro, se identifica la necesidad de reemplazar el historial sintético de visitas
por datos reales acumulados de usuarios (con consentimiento explícito), lo que permitiría que
las reglas de Apriori reflejen patrones verdaderos de comportamiento turístico en Chiapas.
También sería valioso extender el grafo de municipios a rutas multimodales, combinando
segmentos en automóvil con tramos en transporte colectivo para turistas sin vehículo propio.

% ──────────────────────────────────────────────────────────
%  REFERENCIAS
% ──────────────────────────────────────────────────────────
\begin{thebibliography}{9}

\bibitem{ricci2015}
Ricci, F., Rokach, L., \& Shapira, B. (2015).
\textit{Recommender Systems Handbook} (2nd ed.).
Springer. \url{https://doi.org/10.1007/978-1-4899-7637-6}

\bibitem{macqueen1967}
MacQueen, J. (1967).
Some methods for classification and analysis of multivariate observations.
\textit{Proceedings of the Fifth Berkeley Symposium on Mathematical Statistics and
Probability}, 1, 281–297.

\bibitem{agrawal1994}
Agrawal, R., \& Srikant, R. (1994).
Fast algorithms for mining association rules.
\textit{Proceedings of the 20th International Conference on Very Large Data Bases (VLDB)},
487–499.

\bibitem{martello1990}
Martello, S., \& Toth, P. (1990).
\textit{Knapsack Problems: Algorithms and Computer Implementations}.
Wiley.

\bibitem{sectur2023}
Secretaría de Turismo de Chiapas. (2023).
\textit{Estadísticas de afluencia turística 2022–2023}.
Gobierno del Estado de Chiapas.

\bibitem{bi2019}
Bi, J.W., Liu, Y., Fan, Z.P., \& Zhang, J. (2019).
Wisdom of crowds: Conducting importance-performance analysis (IPA) through online reviews.
\textit{Tourism Management}, 70, 460–478.
\url{https://doi.org/10.1016/j.tourman.2018.09.010}

\bibitem{flutter2023}
Google LLC. (2023).
\textit{Flutter: Build apps for any screen}.
\url{https://flutter.dev}

\bibitem{fastapi2023}
Ramírez, S. (2023).
\textit{FastAPI: Modern, fast web framework for building APIs with Python 3.6+}.
\url{https://fastapi.tiangolo.com}

\bibitem{osrm2023}
Project-OSRM. (2023).
\textit{Open Source Routing Machine: High performance routing engine}.
\url{http://project-osrm.org}

\end{thebibliography}

% ──────────────────────────────────────────────────────────
%  CAPÍTULO VI – ANEXOS
% ──────────────────────────────────────────────────────────
\appendix
\chapter{Encuestas de Validación}
\label{anx:encuestas}

Se aplicó una encuesta de satisfacción a los ocho usuarios participantes en la prueba piloto.
El cuestionario incluía cinco preguntas con escala Likert (1–5) sobre: precisión de los
destinos sugeridos, utilidad del desglose de costos, claridad del mapa de ruta, calidad de las
imágenes mostradas y probabilidad de recomendar la aplicación.

% Insertar tabla de resultados de encuesta aquí.

\chapter{Cronograma de Actividades}
\label{anx:cronograma}

% Insertar diagrama de Gantt o captura del tablero Kanban aquí.

\begin{table}[H]
\centering
\caption{Cronograma de sprints}
\label{tab:gantt}
\begin{tabular}{clll}
\toprule
\textbf{Sprint} & \textbf{Período} & \textbf{Objetivo principal} & \textbf{Estado} \\
\midrule
1 & Ago 19 – Sep 1  & Definición de alcances y catálogo inicial & Completado \\
2 & Sep 2 – Sep 15  & Estructura Flutter + navegación base      & Completado \\
3 & Sep 16 – Sep 29 & Backend FastAPI + esquemas Pydantic        & Completado \\
4 & Sep 30 – Oct 13 & Motor ML: K-Means + Apriori               & Completado \\
5 & Oct 14 – Oct 27 & Motor ML: Knapsack + Orquestador          & Completado \\
6 & Oct 28 – Nov 10 & Servicio PLN + integración Groq/Ollama    & Completado \\
7 & Nov 11 – Nov 24 & Grafo de municipios + enrutamiento OSRM   & Completado \\
8 & Nov 25 – Dic 8  & Pruebas, correcciones y despliegue        & En curso   \\
\bottomrule
\end{tabular}
\end{table}

\chapter{Diagramas UML}
\label{anx:uml}

% Insertar diagramas de clases, secuencia y despliegue aquí.

\end{document}
